\chapter{Hinweise für Lehrpersonen}\label{app:lehrpersonen}

\begin{myremark}
	Entwurf vom 24.10.2026 zur Diskussion im Team Modul~B (Cyril Blum, Stephan Pestalozzi).
	Modul~B läuft am Montag zweimal (Morgen: U1c/U1d, Nachmittag: U1a/U1b) mit je zwei parallelen Gruppen;
	beide Lehrpersonen arbeiten mit demselben Ablauf.
\end{myremark}

\section{Konzept}

Der Halbtag ist als \tib{Lernatelier} mit einer durchgehenden Fallgeschichte aufgebaut.
Rund zwei Drittel der Zeit arbeiten die SuS selbstständig in Zweierteams an drei Stationen.
Jede Station schliesst mit einem Eintrag in die \enquote{Fallakte}; am Schluss wird der Fall gelöst und ein kleines Produkt erstellt.
Pflichtaufgaben sichern die Lernziele, Challenges differenzieren nach oben.
Die Lehrperson begleitet, gibt Tipps und beobachtet den Laufzettel.

\section{Ablauf (ca. 180 Minuten netto)}

\begin{center}
	\renewcommand{\arraystretch}{1.35}
	\begin{tabular}{p{1.9cm}p{3.4cm}p{6.4cm}p{2.4cm}}
		\textbf{Zeit}  & \textbf{Phase}           & \textbf{Inhalt}                                                                                                                    & \textbf{Sozialform} \\
		\midrule
		0--20'         & Einstieg                 & Fall Lena vorstellen; Live-Demo Fingerprinting (Cover Your Tracks am Beamer); Ablauf, Regeln, Laufzettel; Zweierteams bilden.       & Plenum              \\
		20--130'       & Lernatelier              & Stationen A, B, C à ca.\ 30--35', gestaffelter Start (je ein Drittel der Teams pro Station); Pause einbauen.                          & Zweierteams         \\
		130--165'      & Fall lösen, Produkt      & Fallakten zusammentragen, Nachricht an Lena; Produkt 1 (Entscheidungsbaum) oder 2 (Phishing-Falle); Phishing-Karten tauschen.       & Zweierteams         \\
		165--180'      & Abschluss                & Kurze Galerie, Auflösung des Falls im Plenum, Exit-Ticket.                                                                          & Plenum              \\
	\end{tabular}
\end{center}

Die Zeiten an die effektiven Lektionen und Pausen des Montags anpassen.

\section{Vorbereitung}

\begin{itemize}
	\item SuS vorab informieren: Handy und Laptop mitbringen, Akku geladen.
	\item Testfoto \enquote{Lena\_Garten.jpg} (gezeichneter Garten mit \acrshort{EXIF}-Daten: Apple iPhone~13, Aufnahme am 30.~Januar 2027 um 14:23, \acrshort{GPS} des öffentlichen Chinagartens Zürich) auf der Lernplattform bereitstellen. Vorher prüfen, ob die Plattform die Metadaten beim Hochladen entfernt; sonst als ZIP-Datei anbieten (\cref{ex:testfoto}).
	\item Erreichbarkeit der Webseiten im Schulnetz testen: coveryourtracks.eff.org, passwortcheck.ch, phishingquiz.withgoogle.com, openverse.org, commons.wikimedia.org.
	\item Material: Skript (ohne Lösungen, Klassenoption \texttt{exerciseonly}), A3-Plakate, Stifte, Karten für Produkt~2, Taschenrechner.
	\item Optional: Tippkarten pro Station (z.\,B. Wege zu den Einstellungen auf iPhone/Android).
	\item Optional: Stationen zusätzlich als Moodle-Kurs mit Selbsttests anbieten (moodle-one).
\end{itemize}

\section{Abgrenzung zu anderen Modulen}

\begin{itemize}
	\item \textbf{Zischtig.ch:} Privatsphäre-Einstellungen in Social-Media-Apps, Sexting, Hassrede, Sucht -- hier bewusst nicht vertieft.
	\item \textbf{Modul~A (digitale Achtsamkeit):} Leitfaden zur eigenen Mediennutzung; deshalb hier andere Produkte.
	\item \textbf{Modul~C (Mechanismen und Beeinflussung):} Deepfakes, Fake-Profile, Bots -- Station~C behandelt nur die Rechtsfragen.
	\item \textbf{Fach Informatik:} Verschlüsselung und Netzwerksicherheit (Kryptologie) -- hier nur Passwörter und \gls{TFA} auf Anwendungsebene.
\end{itemize}

\section{Lehrplanbezug}

\begin{itemize}
	\item Lehrplan~21, Medien und Informatik: \enquote{können mit eigenen und fremden Inhalten Medien\-bei\-träge herstellen und berücksichtigen dabei die rechtlichen Rahmenbedingungen sowie Sicher\-heits- und Ver\-hal\-tens\-re\-geln.}
	\item Rahmenlehrplan Gymnasiale Ma\-tu\-ri\-täts\-schu\-len (2024), Digitalität: \enquote{Digitale Identität reflektieren} (Datenspuren, Privatsphäre, informationelle Selbstbestimmung) und \enquote{Informationsgesellschaft verstehen} (Datensicherheit).
\end{itemize}

\section{Fachliche Hinweise}

\begin{itemize}
	\item Die Zeiten in \cref{ex:passwort_rechnen} beruhen auf der Annahme von $10^{10}$ Versuchen pro Sekunde gegen einen gestohlenen Passwortspeicher (Offline-Angriff). Sie dienen dem Grössenvergleich, nicht als exakte Werte.
	\item Rechtliche Aussagen in Station~C sind bewusst vereinfacht (\gls{URG}, Persönlichkeitsschutz nach \gls{ZGB}, \gls{DSG}). Grenzfälle (Parodie, Zitat, KI-Bilder) im Gespräch differenzieren.
	\item Hilfsangebot für SuS: Pro Juventute, Telefon und Chat 147.
\end{itemize}

\section{Hintergrundlektüre für Lehrpersonen}

Aus der Artikelsammlung (nicht im Repository, nur für den internen Gebrauch):
\begin{itemize}
	\item NZZ am Sonntag, 18.10.2026: Fingerprinting (Tracking ohne Cookies).
	\item NZZ am Sonntag, 18.10.2026: Hacking mit KI (KI-Phishing, Zahlen des BACS).
	\item NZZ, 03.12.2025: Geheime Daten landen auf US-Servern (Schatten-KI).
	\item NZZ, 16.11.2024: Temu -- die Shopping-App, die alle Regeln bricht (Dark Patterns).
\end{itemize}
